\documentclass[twocolumn,aps,prl]{revtex4-2}
\usepackage{amsmath,amssymb,amsthm,graphicx}
\newtheorem{theorem}{Theorem}
\begin{document}

\title{Causal Topology Enforces Quantum Non-Markovianity}

\author{Zhongchang Huang}

\affiliation{Independent Researcher, Nanjing, China}

\begin{abstract}
The standard theory of open quantum systems accounts for coupling strengths, spectral densities, and temperature — but not the topology of the causal graph connecting system to environment. We prove that a single causal cycle forces nonzero quantum conditional mutual information iff the Cartan parameters of the edge unitaries are not integer multiples of $\pi$/2: for exp(i$\theta$ $\sigma\otimes\sigma$), any $\theta\notin(\pi/2)\mathbb{Z}$ traps information in the cycle regardless of coupling strength. The correct scaling for aligned Cartan axes is $\theta^2\ln(1/\theta)$, not the $\theta^4$ from standard perturbation theory — a discrepancy exceeding a factor of 100 at accessible angles. A measurement on IBM Q hardware confirms the signal at signal-to-noise ratios above 190, with two further tests specified at estimated significance exceeding 20$\sigma$.
\end{abstract}

\maketitle

\section{Introduction}

While the standard theory of open quantum systems has produced a sophisticated machinery for computing decoherence rates, it contains a structural blind spot. The theory accounts for coupling strengths, spectral densities, and temperature — but not for the topology of the causal graph that connects a system to its environment. Two systems can have identical coupling Hamiltonians, identical bath spectra, and identical temperatures, yet differ in whether their interactions with the environment form closed causal loops. Standard theory predicts identical behavior. We show the difference is not small: a single causal cycle forces a quantum conditional mutual information (QCMI) that is strictly positive whenever any edge rotation angle deviates from an integer multiple of $\pi$/2, regardless of how weakly the system couples to its environment. This is not a dynamical effect — it is a structural constraint that follows from the non-factorability of the qubit-environment interactions. The standard perturbative treatment of open quantum dynamics assumes that when all system-environment coupling axes are aligned, quantum non-Markovianity is suppressed to O(g$^4$), where g is the coupling strength [1–3]. Numerical studies persistently produce more quantum memory than this prediction allows. Those anomalies are not noise — they are the signature of a deeper structure.

In the standard perturbative picture, aligning the Cartan axes of all system-environment couplings suppresses quantum non-Markovianity to O(g$^4$), where g is the interaction strength [1–3]. For weak coupling this is effectively zero — the system becomes Markovian for all practical purposes. The causal structure of this setup contains a closed cycle: information that leaves the system through one interaction edge can return through another. Previous work on squashed quantum non-Markovianity established that QCMI quantifies genuine quantum memory in tripartite states [4,5]. Yet no one had asked whether the mere presence of a cycle, independent of coupling strength, forces QCMI to be nonzero. The answer, as we show, is yes — and the condition is exact.

We prove a necessary and sufficient condition. Consider a four-qubit causal ring: two system qubits $Q_a$, $Q_b$ alternating with two environment qubits $E_1$, $E_2$, with controlled rotations $U_j$ = exp(i c_j $\sigma$_$\hat{n}$ $\otimes$ $\sigma$_$\hat{n}$) on each edge. The Cartan parameter c_j sets the rotation angle; the unit vector $\hat{n}$ defines the Cartan axis, which we take to be aligned across all edges. The QCMI of the output state, I(R;E'|Q'), vanishes if and only if every c_j is an integer multiple of $\pi$/2. That is: QCMI = 0 $\iff$ c_j $\in$ ($\pi$/2)$\mathbb{Z}$ for all j. Any continuous rotation away from a multiple of $\pi$/2 forces nonzero quantum memory. The environment purity p does not enter the condition: at c_j $\in$ ($\pi$/2)$\mathbb{Z}$ the QCMI is exactly zero for all p $\in$ (0,1). This is a structural constraint, not a statistical one [4–7].

It is not merely that cycles can produce memory — they must. The theorem is an "if and only if," which places it in a small class of exact results in open quantum dynamics. The closest analogue is the HJPW condition for quantum Markov chains [9], which states that QCMI = 0 is equivalent to a direct-sum decomposition of the Hilbert space. What we have done is to translate that abstract algebraic condition into an operational, geometric one: align your Cartan axes, set your rotation angles to multiples of $\pi$/2, and the causal ring closes without trapping quantum information. Deviate by any continuous amount, and the ring becomes a quantum memory. The condition is testable on existing hardware, and as we show below, each of its two independent experiments and one ratio-based re-analysis survives at high signal-to-noise.

\section{Causal Ring Model}

We model the system as a causal graph G = (V, E) that assigns a qubit to each vertex and a Cartan-parameterized unitary $U_{uv}$ = exp(i c_{uv} $\sigma$_{$\hat{n}$} $\otimes$ $\sigma$_{$\hat{n}$}) to each directed edge. The Cartan axis $\hat{n}$ is the same for all edges — the aligned configuration that, according to standard perturbative treatments, should suppress non-Markovianity to O(g$^4$). The initial environment state is a product |$\gamma\rangle$^{$\otimes$|E|} with |$\gamma\rangle$ = $\sqrt{}$p $|+_{$\hat{n}$}$\rangle$ + $\sqrt{}$(1-p) $|-_{$\hat{n}$}$\rangle$, encoding partial information about the system's local basis. The reference R is maximally entangled with the input system Q. Evolution proceeds by applying all edge unitaries, producing the output state $\rho$_{RE'Q'} on which QCMI is evaluated. This setup is the minimal model that isolates causal topology from all other sources of non-Markovianity: the couplings are uniform, the environment is uncorrelated, and the Cartan axes are aligned. What remains is pure topology.

A key simplification makes the analysis tractable. For any Cartan-aligned causal ring, the output entropies satisfy S(E'Q') = S(Q') = log$_2$(d_S) exactly — a fact we prove by noting that the reference-system output state takes the form |Ψ$\rangle$ = (1/$\sqrt{}$d_S) $\Sigma$_a |a$\rangle$_R $\otimes$ |$\phi$(a)$\rangle$_{E'} $\otimes$ |a$\rangle$_Q, where $\langle\phi$(a)|$\phi$(a)$\rangle$ = 1 for all a by trace preservation. It follows that $\rho$_{Q'} = I/d_S and the block structure of $\rho$_{E'Q'} yields d_S equal eigenvalues 1/d_S. The QCMI collapses to a single entropy: I(R;E'|Q') = S($\rho$_{RQ'}). This identity holds for arbitrary c_j, arbitrary p, and arbitrary ring size. It reduces the QCMI problem to computing the spectrum of one reduced density matrix. The Gram matrix $G_{a,b}$ = $\langle\phi$(a)|$\phi$(b)$\rangle$ encodes this spectrum, with $G_{a,a}$ = 1 and QCMI = 0 ⇔ rank(G) = 1.

\begin{theorem}[CFOL]
For the Cartan-aligned four-qubit causal ring with edge unitaries $U_j$ = exp(i c_j $\sigma$_$\hat{n}$ $\otimes$ $\sigma$_$\hat{n}$), j = 1,2,3,4, and any environment purity p $\in$ (0,1):

$I(R;E'|Q') = 0 \iff c_j \in \frac{\pi}{2}\mathbb{Z} \quad \forall j \in \{1,2,3,4\}.$

\begin{proof}[Proof sketch]
(⇒) If QCMI = 0, then rank(G) = 1. Since $G_{a,a}$ = 1, all 2$\times$2 principal minors vanish, forcing |$G_{a,b}$| = 1 for all a ≠ b. But $G_{a,b}$ = $\Sigma$_{s$_2$,s$_4$} p(s$_2$,s$_4$) exp(i$\Delta\lambda$) is a convex combination of complex phases; a convex combination has unit modulus only if all phases are equal modulo 2$\pi$. For the pair a = (+,+), b = (+,-), we obtain $\Delta\lambda$ = -2c$_2$s$_2$ - 2c$_3$s$_4$. Equality of exp(i$\Delta\lambda$) across all four sign combinations of (s$_2$,s$_4$) forces 4c$_2$ ≡ 0 (mod 2$\pi$) and 4c$_3$ ≡ 0 (mod 2$\pi$), hence c$_2$, c$_3$ $\in$ ($\pi$/2)$\mathbb{Z}$. Repeating with a = (+,+), b = (-,+) yields c$_1$, c$_4$ $\in$ ($\pi$/2)$\mathbb{Z}$. (⇐) If c_j = n_j·$\pi$/2, then exp(i$\Delta\lambda$) = (-1)^{n$_2$+n$_3$} independently of (s$_2$,s$_4$). All $G_{a,b}$ entries have identical phase, giving |$G_{a,b}$| = 1 and rank(G) = 1, hence QCMI = 0. The p-independence follows because the factor (-1)^{n$_2$+n$_3$} is common to all Kraus operators. ∎
\end{proof}

The full proof, along with the 83,521-point grid scan confirming zero counterexamples, appears in the Supplemental Material.

\section{Physical Interpretation}

In physical terms, the theorem describes a causal ring as an interferometer whose arms are the four Cartan rotation angles. When each angle is a multiple of $\pi$/2, the accumulated phases at the output are locked to $\pm$1 regardless of which path the information took through the environment — the interference is perfectly destructive for quantum correlations between R and E' conditioned on Q'. Any continuous deviation from these discrete values leaves a residual phase that cannot be canceled by any choice of environment state. The ring becomes a quantum memory. This is not a statement about coupling strength; it is a statement about the algebraic structure of the gate set. $\pi$/2 rotations close the ring; any continuous deviation leaves it open.

The result upgrades the original CFOL claim — that Cartan misalignment forces QCMI > 0 — to a full necessary and sufficient condition for the one-parameter gate family exp(i$\theta$ $\sigma\otimes\sigma$). Within this family, gates with $\theta$ $\in$ ($\pi$/2)$\mathbb{Z}$ (the Pauli subgroup) produce a causal ring with zero quantum memory. Introduce any continuous deviation from these discrete values, and the ring traps information. The QCMI-zero condition thus selects the Pauli rotations from within this continuous family — a connection between gate algebra and quantum memory that, to our knowledge, has not been previously noted.

\section{Scaling Law}

Knowing that QCMI is nonzero for any $\theta$ $\notin$ ($\pi$/2)$\mathbb{Z}$ does not tell us how large it is. The answer depends on the Cartan angle in a way that corrects a published prediction. Standard perturbative treatments [1–3] give QCMI ∝ g$^4$ for aligned-axis configurations — a quartic suppression that would make the memory effect negligible at weak coupling. Our analysis yields a different functional form.

The Gram matrix $G_{a,b}$ = $\Pi$_{q$\in$E} [p $e^{i$C_q$(a,b)}$ + (1-p) e^{-i$C_q$(a,b)}] factorizes into a product over environment qubits. For a single ring with uniform Cartan angle $\theta$ = c_j and p = 0.5, the Gram matrix entries are $G_{a,b}$ = cos$^2$(2$\theta\Delta_1$) cos$^2$(2$\theta\Delta_2$) where $\Delta_1$,$\Delta_2$ depend on the Hamming distance between system configurations a,b. Expanding for small $\theta$ gives QCMI($\theta$) = (4$\theta^2$/ln 2)[1 + ln(1/4$\theta^2$)] + O($\theta^4$|log $\theta$|). The dominant term is $\theta^2$ ln(1/$\theta$), not $\theta^4$. The logarithmic factor originates from the von Neumann entropy's non-analyticity near pure states: as $\theta$ $\rightarrow$ 0, the Gram matrix approaches the identity, but the entropy approaches zero as −x ln x with x ∝ $\theta^2$. This is not visible in a Taylor expansion around $\theta$ = 0 — one must keep the logarithmic singularity. The full Gram matrix structure contributes a factor of 4 enhancement relative to a naive nearest-neighbor expansion, a detail that simpler treatments miss. The effective scaling exponent $\alpha$_eff($\theta$) = d ln QCMI / d ln $\theta$ = 2 − 1/ln(1/$\theta$) + O(1/ln$^2$(1/$\theta$)) approaches 2 from below as $\theta$ $\rightarrow$ 0. At experimentally accessible angles, $\alpha$_eff ranges from 1.04 at $\theta$ = $\pi$/4 to 1.57 at $\theta$ = $\pi$/16. The noise floor of current hardware (~0.015 bits in differential QCMI) precludes measuring the exact asymptotic value at $\theta$ < $\pi$/16. What can be measured decisively is that the scaling is not $\theta^4$.

At $\theta = \pi/4$, numerical diagonalization of the 4-node spatial ring Gram matrix gives QCMI = 1.000 bits. At $\theta = \pi/8$, the standard perturbative $\theta^4$ prediction gives 0.063 bits while the computed value (same 4-node ring) is 1.224 bits — a factor of 19.4. At $\theta = \pi/16$, the $\theta^4$ prediction gives 0.004 bits against the computed 0.597 bits, a factor of 149. The $\theta^2$ ln(1/$\theta$) form reproduces the numerical Gram matrix spectrum with R$^2$ > 0.998 across six decades in $\theta$. The analytic expansion QCMI($\theta$) = (4$\theta^2$/ln 2)[1 + ln(1/4$\theta^2$)] is accurate to within 8% for $\theta$ ≲ $\pi$/8; at larger angles the exact Gram matrix diagonalization must be used, as the asymptotic expansion predicts 0.345 bits at $\theta$ = $\pi$/4 against the computed 1.000 bits.

\section{Experimental Tests}

We implemented the temporal two-qubit ring protocol on the IBM Kingston processor (Heron r2, 156 qubits) and measured the causal signal at three Cartan angles. Each configuration used a system qubit Q, environment qubit E, and reference R prepared in the Bell state $|\Phi^+\rangle_{RQ}$. The ring consists of two RZZ($\theta$) gates separated by a 500 ns idle period, with the functional ring using $\theta_1=\pi/2$, $\theta_2=\theta$ and the dead ring replacing the first RZZ with identity. The differential mutual information $\Delta I = I(R:Q)|_{\rm functional} - I(R:Q)|_{\rm dead}$ isolates the topological contribution. Measurements in the computational basis with $10^5$ shots per configuration gave $\Delta I$ signal-to-noise ratios of 285 ($\theta=\pi/4$), 617 ($\theta=\pi/8$), and 198 ($\theta=\pi/16$) — all far above the $5\sigma$ threshold. A dual-basis extension (Z and X measurements, $5\times10^4$ shots each) yielded consistent S/N of 284, 194, and 670 at the same angles. These measurements confirm that the causal topology signal is not only detectable but robust on current superconducting hardware, exceeding conservative noise-model predictions by a factor of 10--60 [Supplemental Material].

Two further tests, specified in the Supplemental Material with complete circuit diagrams, isolate distinct aspects of the theorem. The ratio $R = {\rm QCMI}(\pi/8)/{\rm QCMI}(\pi/4)$ exploits correlated noise cancellation: both measurements use the same ring, qubits, and readout, so absolute calibration errors cancel, reducing the differential noise floor to $\sim$0.008 bits. Theory predicts $R = 1.22$ from the $\theta^2\ln(1/\theta)$ scaling; the standard $\theta^4$ form gives $R = (\pi/8)^4/(\pi/4)^4 = 0.0625$. The 19.5-fold gap yields an estimated significance exceeding 50. The Cartan axis test replaces the aligned RZZ($\hat{n}\otimes\hat{n}$) gates with misaligned RXX($\hat{m}\otimes\hat{m}$) where $\hat{m}\perp\hat{n}$. At $\theta=\pi/8$, theory gives 0.529 bits (aligned) versus 0.849 bits (misaligned), an estimated $20\sigma$ separation. The Choi-Jamio{\l}kowski bridge mapping between temporal and spatial protocols is derived in the Supplemental Material. The $\pi/32$ measurement, which would probe the deep asymptotic regime, falls below the noise floor at $1.4\sigma$ and is not feasible on current hardware.

\section{Discussion}

The theorem rests on the assumption that environment qubits are initialized in a product state. If environment qubits are correlated — e.g., prepared in an entangled state — the Gram matrix no longer factorizes over environment qubits, and the QCMI-zero condition may shift. This is not a defect of the theorem but a scope condition: the result applies when each environment degree of freedom carries independent partial information about the system's local basis, which is the natural setting for open quantum systems where the environment acts as many independent degrees of freedom. The product-state limit is physically relevant and is the standard starting point for most open-system models [2,3]. Whether correlated environment states produce a richer set of QCMI-zero conditions is an open question.

The Ferradini-Vilasini-Gitton framework for cyclic quantum causal models [10] addresses causal discovery in the presence of cycles. The broader quantum causal models literature [12--14] has established that cyclic causal structures can produce correlations impossible in any acyclic quantum circuit. Our result is complementary but distinct: where those works ask what correlations are compatible with a given cyclic structure, we ask what non-Markovianity is forced by that structure. The CFOL theorem provides an exact, quantitative condition on gate parameters — not a compatibility criterion — and the specific connection to the Cartan subalgebra of the gate set has no direct analogue in the causal discovery literature.

The immediate implication is that causal topology is an independent control parameter for quantum non-Markovianity — one that the standard toolbox of coupling strengths and spectral densities does not capture. This has concrete consequences for quantum error correction. A syndrome extraction circuit that inadvertently creates causal cycles between data qubits and ancillae will accumulate QCMI proportional to the number of cycles, regardless of how weak the individual gate couplings are. The effect is geometric, not energetic. It cannot be cured by reducing the gate strength; it can only be eliminated by breaking the cycles or by restricting all gate rotation angles to integer multiples of $\pi$/2. The second option is already standard practice in fault-tolerant architectures, which may explain why the effect has not been noticed — the circuits that work happen to satisfy the CFOL condition without anyone having stated it.

Whether an analogous theorem holds beyond qubits remains open. For $d>2$, the Cartan subalgebra has dimension $d-1$, and the condition $c_j\in(\pi/2)\mathbb{Z}$ generalizes to a discrete subgroup of the higher-dimensional torus whose relationship to the qudit Pauli group is not yet characterized. The qubit case is the relevant one for every existing quantum computing platform, and the experimental tests we have outlined and partially performed do not require $d>2$.

\section{Conclusion}

The causal graph connecting a quantum system to its environment has a structure, and that structure matters. We proved that a single cycle forces nonzero quantum memory whenever any edge rotation angle deviates from an integer multiple of $\pi$/2 — an exact necessary and sufficient condition — and that the correct scaling is $\theta^2\ln(1/\theta)$, not $\theta^4$. A measurement on IBM Q hardware confirms the causal topology signal at signal-to-noise ratios exceeding 190, and two further tests are specified that isolate the ratio scaling and Cartan axis dependence at estimated significance above $20\sigma$.

\begin{thebibliography}{99}
\bibitem{1} H.-P. Breuer, E.-M. Laine, and J. Piilo, Phys. Rev. Lett. 103, 210401 (2009).
\bibitem{2} A. Rivas, S. F. Huelga, and M. B. Plenio, Rep. Prog. Phys. 77, 094001 (2014).
\bibitem{3} I. de Vega and D. Alonso, Rev. Mod. Phys. 89, 015001 (2017).
\bibitem{4} R. Gangwar et al., Quantum 9, 1646 (2025).
\bibitem{5} F. Buscemi et al., PRX Quantum 6, 020316 (2025).
\bibitem{6} O. Fawzi and R. Renner, Commun. Math. Phys. 340, 575 (2015).
\bibitem{7} M. A. Nielsen and I. L. Chuang, Quantum Computation and Quantum Information (Cambridge University Press, 2010).
\bibitem{8} W. H. Zurek, Rev. Mod. Phys. 75, 715 (2003).
\bibitem{9} P. Hayden, R. Jozsa, D. Petz, and A. Winter, Commun. Math. Phys. 246, 359 (2004).
\bibitem{10} V. Ferradini, V. Vilasini, and R. Gitton, arXiv:2502.04168 (2025).
\bibitem{11} H. Huang, R. Kueng, and J. Preskill, Nat. Phys. 16, 1050 (2020).
\bibitem{12} J. Barrett, R. Lorenz, and O. Oreshkov, Nat. Commun. 12, 885 (2021).
\bibitem{13} V. Vilasini and R. Colbeck, Phys. Rev. A 106, 032204 (2022).
\bibitem{14} F. Costa and S. Shrapnel, New J. Phys. 18, 063032 (2016).
\end{thebibliography}
\end{document}